\section{Hilbert multipliers and the absence of critical points}\label{sec:conclusion}

We now use the conditional negativity proved in \Cref{thm:negative-kernel} to construct all the scalar multipliers needed in \Cref{prop:multiplier} at once. The Hilbert-space construction is the classical squared-distance correspondence of Schoenberg \cite{Schoenberg1938}; we give the short Gram-kernel proof.

\begin{lemma}\label{lem:hilbert}
For every $p\in\D$ there is an injective Lipschitz map $b:\Sone\to\HH$, where $\HH$ is a separable real Hilbert space, such that
\begin{equation}\label{eq:hilbert-distance}
\norm{b(s)-b(t)}_\HH^2=D_p(s,t).
\end{equation}
\end{lemma}
\begin{proof}
Fix $s_0\in\Sone$. Conditional negativity makes the anchored kernel
\[
\Gamma(s,t)=\tfrac12\bigl(D_p(s,s_0)+D_p(t,s_0)-D_p(s,t)\bigr)
\]
positive semidefinite: append to any finite family of coefficients the coefficient at $s_0$ that makes their sum zero. Give formal finite linear combinations of symbols $b(s)$ the bilinear form with Gram kernel $\Gamma$, quotient by its nullspace, and complete. Since $D_p(s,s)=0$, this realizes \eqref{eq:hilbert-distance} and $b(s_0)=0$.

Positivity of $K$ and finiteness and positivity of $N$ off the diagonal give $D_p(s,t)>0$ for $s\ne t$, so $b$ is injective. Also $N\ge N_0$ gives
\[
\norm{b(s)-b(t)}_\HH^2
\le\pi\norm{K(p,\cdot)}_\infty^2|s-t|^2.
\]
Thus $b$ is Lipschitz. Replace $\HH$ by the closed span of $b(\Sone)$; it is separable by continuity and separability of the circle.
\end{proof}

\begin{proof}[Proof of \Cref{thm:main}]
Take an arbitrary $0\ne u\in\V$, with $g$ defined by \eqref{eq:g}. Suppose $\Phi(p)$ is a critical point. Choose $b$ from \Cref{lem:hilbert}, and let $b_j$ be its coordinates in a finite or countable orthonormal basis. Each $b_j$ is real Lipschitz, so \Cref{prop:multiplier} applies. For partial sums, the absolute integrand is bounded by
\begin{align*}
N(s,t)\sum_{j\le m}|b_j(s)-b_j(t)|^2|g(s)||g(t)|
&\le N(s,t)D_p(s,t)|g(s)||g(t)|\\
&=K(p,s)K(p,t)|g(s)||g(t)|.
\end{align*}
The right side is bounded and integrable. Dominated convergence therefore gives
\begin{align}
\sum_j E(b_jg)
&=\frac12\iint K(p,s)K(p,t)g(s)\cdot g(t)\dd s\dd t\notag\\
&=\frac12\left|\int_{\Sone}K(p,s)g(s)\dd s\right|^2
=\frac12|W_g(p)|^2=0.\label{eq:sum-energy}
\end{align}
Every summand is nonnegative. By \eqref{eq:E-nullspace}, for every $j$ there is $v_j\in\V$ whose boundary gradient, pulled back by $\Phi$, is $b_jg$. The trace equality holds pointwise because both sides are continuous. The boundary gradient of $u$ is $g$.

Set $S_g=\{s\in\Sone:g(s)\ne0\}$. This is a nonempty open set. If it were empty, each Cartesian derivative of $u$ would solve the Helmholtz equation with zero Dirichlet trace, contradicting $\mu<\lambda_D$ unless both vanished. Then $u$ would be a constant eigenfunction for a positive eigenvalue, hence zero.

The two possibilities for $S_g$ lead respectively to multiplicity and boundary-uniqueness contradictions.

\smallskip\noindent
\emph{Case 1: $S_g$ is dense.}
The set $b(S_g)$ cannot lie in an affine line. Otherwise continuity would place $b(\Sone)$ in that closed line, but a continuous injective circle cannot map into a line. Indeed its compact connected image would be an interval; deleting an interior point disconnects an interval but not a circle with one point removed.

Choose three points of $S_g$ whose images are affinely independent. Their two difference vectors are linearly independent in $\HH$, so some two coordinates $j,k$ have a nonzero $2\times2$ minor. To see this, if all such minors vanished, a nonzero coordinate of the first difference vector would force every coordinate of the second to be proportional to it. Thus the functions $1,b_j,b_k$ are linearly independent on those three points. Since $g\ne0$ there, the vector traces $g,b_jg,b_kg$ are linearly independent. They are gradient traces of three functions in $\V$, contradicting \Cref{prop:multiplicity}.

\smallskip\noindent
\emph{Case 2: $S_g$ is not dense.}
There is a nonempty open arc $I$ on which $g=0$. Since $S_g$ is nonempty and open and $b$ is injective, some coordinate $b_j$ is nonconstant on $S_g$. Consequently $g$ and $b_jg$ are linearly independent, and so are their eigenfunctions $u$ and $v_j$. Their gradients vanish on the boundary arc $\Phi(I)$, so each has a constant trace there. Choose a nonzero linear combination $w$ of them whose constant trace on that arc is zero. It also has zero normal derivative there.

Take a small neighborhood of an interior point of the arc that meets no other boundary portion, and extend $w$ by zero to the exterior of $\Omega$ in that neighborhood. This is a local version of the zero-extension argument in \cite[Lemma, p.~414]{Filonov2005}. Integration by parts against compactly supported tests shows that this extension solves $(\Delta+\mu)w=0$ distributionally: both its Dirichlet and normal traces on the intervening boundary vanish. Interior elliptic regularity and analyticity make the extension analytic. Since it is zero on an exterior open set, it is zero throughout a smaller neighborhood. Analytic continuation inside the connected domain gives $w=0$ on $\Omega$, contradicting independence of $u$ and $v_j$. No analyticity of the boundary is used.

Both cases contradict the assumed critical point. Thus $\grad u$ is nowhere zero in $\Omega$. The continuous function $u$ attains both extrema on the compact closure. Any interior extremum would be critical, so both are attained on the boundary. An interior value equal to either boundary extremum would itself be a global extremum. This proves both strict inequalities in \eqref{eq:hotspots}.
\end{proof}
